\section{Sequencing Award-Layer Screening and Bid-Layer Forensics}
\label{sec:screening_forensics}

The previous sections establish two points.
Section~\ref{sec:validation} validates the award-layer ranking
against cartel-adjacent loser-side labels;
Section~\ref{sec:cobidder_type} shows that cobidders are not
merely generic high-volume losers. The remaining question is
institutional. What should an agency do with a signal that is
cheap enough to compute before bid-level evidence is recovered?
The agency's problem is not only prediction; it is evidence
allocation under incomplete observability.

The answer is not to treat the award-layer screen as a
substitute for bid analysis. The answer is to put it before bid
analysis. Award-layer screening creates suspicion on the record
layer the agency already observes; bid-layer forensics then
investigates a smaller survivor pool using richer evidence from
submitted bids. The object is therefore not a detector horse
race under full observability. It is an evidence-allocation
problem under incomplete observability.

This section evaluates the sequence in three steps.
Subsection~\ref{sec:bid_layer_benchmark} asks how much ranking
information is available at the award layer relative to the bid
layer on the same cartel-adjacency target---an
information-cost comparison, not a methodological dominance
test. Subsection~\ref{sec:complementary_signals} asks whether
award-layer and bid-layer signals are complements or
substitutes. Subsections~\ref{sec:sequential_gatekeeper} and
\ref{sec:temporal_gatekeeper} ask whether the award-layer
ranking can gate costly bid recovery while preserving much of
the adjudication signal, both contemporaneously and under
temporal holdout.

The third question is the operational object. The sequential
gatekeeper exercise should be read as a cost-of-evidence
result, not only as a classification table:
Table~\ref{tab:gatekeeper_v18} reports the reduction in
bid-microdata pool size that the award-layer screen produces,
alongside the recall it preserves against the adjudication
record. That trade-off, not the comparison of stand-alone AUCs,
is the institutional question this paper addresses.

\subsection{Same Target, Different Information Costs}
\label{sec:bid_layer_benchmark}

The forensic comparator is the bid-distribution tradition.
Imhof--Wallimann-style detectors use within-tender bid
moments---dispersion, skewness, spread, distance from the
winning bid, and related features---to flag patterns consistent
with coordinated bidding
\citep{imhof2019detecting,wallimann2023machine}. Those features
live at the bid layer. They require per-bid offer values and
enough usable bid histories to construct a firm-level forensic
profile. The award-layer screen uses a thinner envelope: zero
wins and participation counts. The question is not which
detector is more sophisticated. It is how much ranking
information is available before the richer bid layer is opened.

Table~\ref{tab:imhof_cost_v18} puts the two layers on the same
cartel-adjacency target. That common target matters: it
prevents the comparison from changing the object being ranked.
What differs is the data envelope. The bid-layer benchmark
requires per-bid offer values and within-tender moments; the
award-layer measures require only routine award records.
Table~\ref{tab:imhof_cost_v18} is therefore an information-cost
table before it is a classification table.

Within the always-loser pool for which bid-distribution
features can be computed ($N=\valImhofPoolN$), the binary
frequent-loser flag reaches AUC $\valImhofFLBin$ ($95\%$ CI
$\valImhofFLBinCI$). The continuous participation score reaches
AUC $\valImhofFLcont$ ($\valImhofFLcontCI$). The full bid-layer
benchmark reaches AUC $\valImhofFull$ ($\valImhofFullCI$). The
combined award-layer and bid-layer specification reaches AUC
$\valImhofComboCont$ ($\valImhofComboContCI$).

This is not an outperformance claim. The award-layer statistic
does not dominate bid-distribution forensics; the bid-layer
benchmark remains the natural forensic tool when bid microdata
are available. The point is sharper: a record-layer statistic,
computed before forensic recovery of bid values, discriminates
the same adjudication-anchored target in the same range as a
much richer bid-layer benchmark. That information-cost result
is the premise for sequencing. If the award layer is
informative and not redundant with the bid layer---a question
the next subsection addresses---it can decide where bid-layer
forensics should begin.

\begin{table}[htbp]
\centering
\caption{Discrimination at Different Information Costs}
\label{tab:imhof_cost_v18}
\footnotesize
\begin{adjustbox}{max width=\textwidth,center}
\begin{threeparttable}
\begin{tabular}{p{4.6cm}p{5.4cm}cc}
\toprule
Model & Information required & AUC & 95\% CI \\
\midrule
Frequent-loser flag & Award records: participation count and zero wins
& $\valImhofFLBin$ & $\valImhofFLBinCI$ \\
Participation intensity & Award records: $\log(1+\text{tenders\_count})$ and zero wins
& $\valImhofFLcont$ & $\valImhofFLcontCI$ \\
Bid-layer benchmark & Per-bid offer values and within-tender bid moments
& $\valImhofFull$ & $\valImhofFullCI$ \\
Award layer $+$ bid layer & Participation intensity plus bid-distribution moments
& $\valImhofComboCont$ & $\valImhofComboContCI$ \\
\bottomrule
\end{tabular}
\begin{tablenotes}
\footnotesize
\item \textit{Notes:} Five-fold cross-validated AUC against
CADE-cobidder labels within the always-loser pool for which the
bid-distribution benchmark can be computed
($N=\valImhofPoolN$). The combined specification uses the same
target and adds the award-layer score to the bid-layer moments.
\end{tablenotes}
\end{threeparttable}
\end{adjustbox}
\end{table}

\subsection{Complementarity, Not Substitution}
\label{sec:complementary_signals}

Comparable AUCs do not by themselves justify a two-stage
architecture. If the award-layer statistic were merely a noisy
proxy for bid-distribution moments, it would be useful mainly
where bid microdata are unavailable. The stronger test is
whether the award layer adds information once the bid layer is
observed.

It does. Adding the award-layer statistic to the bid-layer
benchmark raises discrimination to $\valImhofComboCont$
($\valImhofComboContCI$). In the identical same-sample audit,
where all models are evaluated only on firms with complete
bid-layer features, the frequent-loser flag adds
$+\valImhofIncFL$ AUC over the full Imhof--Wallimann benchmark
alone (DeLong $p=\valImhofIncFLp$). The combined specifications
gain $\valImhofIncCombo$ AUC over the bid-layer benchmark
($p<0.001$).

The gain is substantively important because the two layers
measure different legal-economic objects. Bid moments describe
conduct inside tenders: dispersion, spread, rank distance, and
other within-tender patterns. Participation intensity describes
repeated loser-side exposure across tenders, before the bid
file is opened. The first is closer to forensic evaluation; the
second is closer to investigative priority. A regulator should
not collapse them. The first creates a reason to look; the
second helps evaluate what is found after looking.

The combined model is therefore a complementarity diagnostic,
not the operational baseline. Running it on every firm requires
full bid-layer observability before triage---precisely the cost
the architecture is designed to manage. The enforcement use of
the same non-redundancy appears in the next subsection:
award-layer screening first creates the survivor pool, and
bid-layer forensics then evaluates that pool.

\subsection{Gatekeeping the Forensic Stage}
\label{sec:sequential_gatekeeper}

The operational test is not whether the cheap screen beats the
forensic benchmark. The relevant question is whether the cheap
screen can reduce the forensic workload without discarding too
much adjudication-anchored signal.
Table~\ref{tab:gatekeeper_v18} should therefore be read as a
cost-of-evidence table before it is read as a classification
table: each rule is evaluated not only by precision and recall,
but also by how many firms require bid microdata.

The baselines differ in observability. The award-layer-only
rule uses no bid microdata. The bid-layer-only rule and joint
scoring require bid microdata for the full evaluation pool.
Joint scoring is an upper bound under full observability: it
gives the classifier both award-layer and bid-layer features
for every firm before triage has occurred. The sequential rule
imposes the agency's institutional constraint: it allocates
bid-layer evaluation to a survivor pool selected by the award
layer.

The sequential rule operates in two stages. First, it ranks
firms using the award layer alone and keeps the top
$K_1=\valArchSeqKone$. Second, it recovers or processes bid
microdata only for those survivor firms, applies the bid-layer
benchmark within the survivor pool, and selects the final
top-$k$ flags. The architecture therefore reserves the costly
forensic step for the smaller set of firms the award layer
already prioritizes; it does not assume bid microdata are
available before triage.

The evaluation pool contains $\valArchPoolSize$ always-loser
firms with computable bid-distribution features and
$\valArchNPos$ CADE-cobidder positives. At the top $1{,}000$
flags, joint scoring finds $\valArchTPThouJoint$ cobidders, or
recall $\valArchRecallThouJoint$, but requires bid microdata
for all $\valArchPoolSize$ firms. The sequential rule finds
$\valArchTPThouSeqK$ of $\valArchNPos$ cobidders, or recall
$\valArchRecallThouSeqK$, while requiring bid microdata for
only $\valArchSeqKone$ firms. Relative to joint scoring, the
recall loss is $\valArchRecallCost$, while the bid-microdata
footprint falls by $\valArchFootprintReduction$.
Table~\ref{tab:gatekeeper_v18} reports the same trade-off at
the top-$500$ cutoff as well.

\begin{table}[htbp]
\centering
\caption{Gatekeeping the Bid-Layer Forensic Stage}
\label{tab:gatekeeper_v18}
\footnotesize
\begin{adjustbox}{max width=\textwidth,center}
\begin{threeparttable}
\begin{tabular}{p{5.2cm}rrrrr}
\toprule
Rule & $k$ & TP & Precision & Recall & Bid-microdata firms \\
\midrule
Award-layer only & $500$ & $66$ & $\valArchPrecAtFiveHundFL$ & $0.342$ & $0$ \\
Bid-layer only & $500$ & $73$ & $\valArchPrecAtFiveHundImhof$ & $0.378$ & $\valArchPoolSize$ \\
Joint scoring, award $+$ bid & $500$ & $112$ & $\valArchPrecAtFiveHundJoint$ & $0.580$ & $\valArchPoolSize$ \\
Sequential: award $\rightarrow$ bid, $K_1=\valArchSeqKone$
& $500$ & $96$ & $\valArchPrecAtFiveHundSeqK$ & $0.497$ & $\valArchSeqKone$ \\
\midrule
Joint scoring, award $+$ bid & $1{,}000$ & $\valArchTPThouJoint$
& $0.142$ & $\valArchRecallThouJoint$ & $\valArchPoolSize$ \\
Sequential: award $\rightarrow$ bid, $K_1=\valArchSeqKone$
& $1{,}000$ & $\valArchTPThouSeqK$
& $0.131$ & $\valArchRecallThouSeqK$ & $\valArchSeqKone$ \\
\bottomrule
\end{tabular}
\begin{tablenotes}
\footnotesize
\item \textit{Notes:} Evaluation pool:
$\valArchPoolSize$ always-loser firms with computable
bid-distribution features; positives are $\valArchNPos$
CADE-cobidders. The bid-microdata column reports how many firms
require per-bid offer values before the rule can be applied.
\end{tablenotes}
\end{threeparttable}
\end{adjustbox}
\end{table}

The sequential rule is not the precision champion---joint
scoring has the easier job because it uses every information
layer for every firm---and it is not claimed to be universally
optimal. It earns its place by operationalizing the paper's
evidence-allocation problem: when bid recovery is costly, a
modest loss in recall can be worth a large reduction in
forensic-data exposure. That is the trade-off faced by an
enforcement agency that cannot inspect every bidder, every bid,
and every tender with the same forensic intensity.

\subsection{Temporal Holdout Gatekeeping}
\label{sec:temporal_gatekeeper}

The in-sample gatekeeper exercise still deserves a timing
check. A classifier trained and evaluated inside the full
sample may rank firms well because it observes long histories
close to the adjudicated outcome. A timing audit therefore asks
whether the architecture still works when features are formed
before the target window. This is not a real-time field
deployment, but it imposes a stricter timing discipline on the
evidence-allocation exercise.

We restrict feature construction to $2009$--$2016$ and evaluate
against the same post-window adjudication-anchored cobidder
target. The temporal-holdout pool contains
$\valArchTHPoolSize$ always-loser firms with complete
train-window bid features and $\valArchTHNPos$ positives. The
target itself remains adjudication-based and observed later;
the discipline is on the score, not on the labels.

The audit preserves the architecture. At top $500$, the
award-only rule loses $\valArchTHFLDrop$ of its in-sample true
positives, joint scoring loses $\valArchTHJointDrop$, and the
sequential rule loses only $\valArchTHSeqDrop$. In levels, the
sequential rule recovers $\valArchTHTPFiveHundSeqK$ positives
at top $500$, compared with $\valArchTHTPFiveHundJoint$ for
joint scoring. At top $1{,}000$, it recovers
$\valArchTHTPThouSeqK$ positives, compared with
$\valArchTHTPThouJoint$ under joint scoring.

This should not be read as proof that sequential gatekeeping
dominates joint scoring. The safer interpretation is narrower
and more useful: conditioning bid-layer forensics on an earlier
award-layer screen remains competitive out of time while
preserving a smaller forensic footprint. The result reinforces
the paper's sequencing claim under timing discipline; it does
not establish universal optimality, real-time enforcement
success, or causal prediction.

The legal interpretation follows the same boundary as the
empirical design. A frequent-loser flag cannot establish an
agreement, identify a damages base, or prove that a particular
bid was cover. It can identify where the agency should spend
the next unit of investigative effort. The bid-layer stage then
examines whether the bidding record contains the richer
patterns needed for proof. For a journal concerned with how law
works under institutional constraint, that sequencing is the
contribution: a cheap administrative record can improve the
allocation of costly legal evidence gathering without
pretending to be legal evidence itself.
